%Previous studies summary
%safety
%static field
From a safety standpoint, the static magnetic field ($B_0$) represents the main hazard and must be shielded properly. A vast majority of diagnostic clinical scans are completed daily without incidents indicating that magnetic interactions with normal tissues are within the bounds of safety up to the highest fields now in use for \gls{mri} \cite{Schenck2000}.  However, there have been well-documented incidents associated with ferromagnetic objects either implanted in the patient or inappropriately brought into the scanner environment (Zone IV) \cite{Delnoa2019,Sawyer2024}. Recent cases continue to underscore that such accidents most commonly result from failures in pre-entry screening and access control protocols, rather than from the technology itself \cite{CNN2025McAllister}.

\gls{mri} safety protocols are grounded in a system of absolute and relative contraindications. Absolute contraindications include non–MR-conditional cardiac implantable electronic devices (\glspl{cied}), metallic intraocular, and foreign bodies, among others. Relative contraindications consider implants such as stents, surgical clips, and programmable shunts that may be scanned under controlled conditions \cite{Ghadimi2025}. The safety classification of each device is guided by testing standards and certification data, with CIEDs representing a critical category due to their susceptibility to electromagnetic interactions \cite{Schenck2000, fda2023testing}. 

The American College of Radiology (ACR) Manual on MR Safety (2024) \cite{ACR_Manual_on_MR_Safety_2024} outlines the current regulatory and procedural framework for \gls{mri} safety. Recommendations include structured zoned access control and interprofessional collaboration through clearly defined positions (MR Medical Director, MR Safety Officer, and MR Safety Expert) to prevent unscreened individuals or ferromagnetic objects from inadvertently entering the magnet room. Pre-scan screening, implant model and status verification, and continuous patient monitoring are still necessary for this safety infrastructure \cite{Ghadimi2025}.

%What pertinent literature needs to be reviewed to orient the reader and provide sufficient background information for them to understand your study?.

%-------------------------------------------------------------------------------------------------------------------
%-------SCHALLER 2021-------------------------------------------------------------------------------------------
%-------------------------------------------------------------------------------------------------------------------

In addition to current safety guidelines, studies have centered on qualitative findings. For example, a large observational study by Schaller \textit{et al.} (2021) evaluated 200 \gls{mri} examinations in 139 patients with a total of 243 abandoned leads \cite{schaller2021}. Their objective was to assess the safety of 1.5-T \gls{mri} in patients with at least one abandoned \gls{cied} lead and to evaluate any potentially harmful effects on the active \gls{cied} leads. 

Before the scans, safety was a priority. Each case was discussed by radiologists and the ordering healthcare professionals in order to balance the advantages of alternative tests against the dangers associated with \gls{mri}. Next, baseline battery voltage, sensing, pacing thresholds, and lead impedances were recorded. During the \gls{mri}, patients were monitored continuously via electrocardiography telemetry and pulse oximetry and the scans were performed in a normal specific absorption rate (SAR) mode, ensuring they did not exceed 2.0 Watts per kg for the whole body. In some cases involving the brain, a transmit-receive radiofrequency coil was used to avoid the central location of the leads relative to the coil. 

Immediately after the \gls{mri}, the \gls{cied} evaluation was repeated to check for changes in capture threshold, sensing, and lead impedance. The research found a low rate of adverse events (3\%), with no serious complications such as sustained arrhythmias or device resets. The observed events were primarily transient, non-life-threatening changes in lead sensing parameters that resolved spontaneously. Notably, this safety profile held true even for thoracic and cardiac \gls{mri} scans, where the theoretical risk is highest.

%-------------------------------------------------------------------------------------------------------------------
%-------Vuorinen 2021-------------------------------------------------------------------------------------------
%-------------------------------------------------------------------------------------------------------------------


Complementing this, a retrospective study by Vuorinen \textit{et al.} (2021) focused on 26 \gls{mri} scans in 17 patients with epicardial leads \cite{vuorinen2021}. With a similar objective, on a 1.5-T \gls{mri}, researchers distinguished between ``old" epicardial leads (implanted before 2000) and ``modern" leads (implanted in 2000 or later) to compare safety outcomes. 

They followed a dedicated institutional protocol, previously implemented, which included measuring the length of the leads using chest x-rays prior to \gls{mri} scan to determine if the lead reached lengths, of 24 cm, which might increase heating risks. Similarly, researchers documented \gls{cied} generator and lead models and parameters with implantation dates.

The study's key finding was that \gls{mri}s performed on patients with ``modern" epicardial leads had no detectable adverse events. However, one definite and one possible adverse event were noted in patients with ``old" leads, suggesting that lead age and technology may be a factor in the risk profile. 

These observations rely on clinical monitoring and empirical outcomes rather than direct quantification of physical interactions with the static magnetic field. To bridge this gap and enable more precise risk evaluation, standardized quantitative protocols are essential for measuring the magnetically induced forces on such objects.

%the current protocols to assess objects near MRI field
\subsection{Quantification of Magnetically Induced Displacement Force and Torque}
\label{sec:staticfield-forces}

%-------------------------------------------------------------------------------------------------------------------
%-------ferreira 2017 ---------------------Field map----------------------------------------------------------------
%-------------------------------------------------------------------------------------------------------------------


ASTM F2052-21 offers a standardized technique for assessing magnetically induced displacement force on medical equipment in order to set safety standards for such situations. The methodology includes a non-magnetic pendulum apparatus with a protractor that allows for the a suspended subject to be pulled by the magnetic field. The characterization of the static magnetic field and its spacial gradient is required. For this purpose, Ferreira \textit{et al.} (2017), developed a custom-designed fixture with an AlphaLab Vector/Magnitude Gaussmeter (Model VGM) probe to map the spatial gradients inside a Siemens Magnetom Trio \textsuperscript{\texttrademark}3T MRI scanner \cite{ferreira2017}. Measurements of the magnetic field vector components (X, Y, and Z) were collected along the Z-axis in 5cm increments, correctly capturing the static magnetic field and ensuring repeatability and showing a steady spatial gradient. 


%-------------------------------------------------------------------------------------------------------------------
%-------Woods 2021 ------------------DISPLACEMENT AND GRADIENT--------------------------------------------
%-------------------------------------------------------------------------------------------------------------------

To quantify the magnetically induced displacement force, Woods \textit{et al.} (2021), conducted a similar experiment following ASTM F2052-21 on eleven metal alloys commonly used in medical implants \cite{Woods2021_MRIdisplacement}. For deflection force testing, nine specimens (in three nominal mass categories: 2, 20, and 200 g) of each alloy were tested in a Siemens MAGNETOM Trio\textsuperscript{\texttrademark} 3T MR system with the objective to establish evidence for MR conditional device labeling.

The main results regarding magnetically induced deflection force showed that for seven of the eleven tested alloys, the magnetic force was found to be less than the device weight in any current clinical 1.5 T, 3.0 T, and 7.0 T MR system. These safe alloys included various grades of CP titanium, Ti-6Al-7Nb, Ti-6Al-4V ELI, Ti-15Mo, Biodur 108, and Co-28Cr-6M. However, for the other alloys, including 316L stainless steel and 22Cr-13Ni-5Mn, the magnetically induced force was calculated to be greater than the device weight in certain high-field systems, such as 3.0 T and 7.0 T clinical scanners.

%-------------------------------------------------------------------------------------------------------------------
%-------Woods 2021 ------------------TORQUE--------------------------------------------
%-------------------------------------------------------------------------------------------------------------------


Following a series of torque measurements, Woods \textit{et al.} (2021) reported that for medical instruments made from their list of alloys the magnetically induced torque will always be less than the worst case torque produced by gravity in a 3.0 T static magnetic field \cite{Woods2021_MRIdisplacement}. Notably, no measurable movement was observed in any of the torque specimens tested at 3.0 T. However, the ASTM conservative safety criterion requires that magnetically induced torque do not exceed its gravitational equivalent, $\tau < Lmg$. Using this criterion, the authors established an upper bound on torque by measuring the coefficient of static friction between specimens and the test surface, finding that magnetic torque remained below roughly 30\% of the gravitational torque across all tested alloys. This result suggests that under typical clinical \gls{mri} circumstances, torque is not a limiting safety concern for implant-grade versions of these materials. This finding demonstrates that torque effects are secondary to displacement forces. Consequently, device designers can prioritize addressing displacement force effects while being confident that torque is generally well-controlled under typical clinical field strengths.

However, the study was probably done using ASTM F2213-17 ``The Low Friction Surface Method" which estimates magnetically induced torque by placing the device on a low-friction, non-metallic, non-conductive surface near the \gls{mri} isocenter. The device is rotated step-wise, and any spontaneous alignment or rotation toward the static magnetic field is observed. If no movement occurs, the maximum possible torque is estimated from the device’s weight and the surface’s friction coefficient outside the MR environment. ASTM F2213-17 states that, if rotation is observed, more precise tests should be used instead which means this method relies on macroscopic motion and may not be able to detect $mN$ forces. 
%-------------------------------------------------------------------------------------------------------------------
%-------SEO 2021------------------TORQUE--------------------------------------------
%-------------------------------------------------------------------------------------------------------------------

In this regard ASTM F2213-17 define other methodologies, mainly the Torsional Spring or Pulley Method. A study by Seo \textit{et al.} (2021), utilized ``the Pulley Method" for the safety assessment of magnetically induced torques on \gls{aimd} in 1.5 T and 3.0 T magnetic resonance imaging (\gls{mri}) environments \cite{Seo2021767}. They constructed a cylindrical shaped apparatus entirely from acrylic, including moving parts like the bearings, to minimize magnetic field resistance. 

This methodology determines the maximum magnetically induced torque by fixing the \gls{aimd} onto a rotating plate placed at the \gls{mri} scanner isocenter. A thread attached to the plate is connected to a force gauge located far enough to avoid interference with the magnetic field. As the plate rotates to align with field, the thread tension is measured.

Prior to \gls{mri} scanner deployment, the apparatus was validated using calibrated masses applied at known radii to generate controlled torques. Across a range of applied loads, the measured torque agreed with theoretical values to within an average error of approximately 1.5\%. The preliminary experiment was configured to replace the magnetic field's effect with known gravitational forces (weights). Weights ranging from 100–500 g were placed on the rotating plate at the point where the \gls{aimd} would be located, which was 2.9 cm from the center pillar. A force gauge, connected via a thread, was placed 4.9 cm from the center and used to pull the plate in the opposite direction of the weights. The force gauge used was load-cell-based with a resolution of 0.01 N. Each test was performed ten times for each weight (100 g, 200 g, 300 g, 400 g, and 500 g), and measurement was performed only when the equipment was stable. 

Additionally, Seo \textit{et al.} (2021) addressed the friction issue not by eliminating it entirely, but by quantifying and compensating for it mathematically within their final torque calculation \cite{Seo2021767}. The friction of the rotating plate ($F_f$) was determined in a prior experiment. Then, the friction value was explicitly used in the formula for magnetically induced torque ($\tau$) calculations, which follows the principles of their Pulley Method. They explained that the torque ($\tau$) is obtained by multiplying the distance from the center of rotation ($r$) to the \gls{aimd}’s position by the difference between the measured force ($F_m$) and the friction of the rotating plate ($F_f$), $\tau = r (F_m - F_f)$.

In the clinic, they tested two types \gls{aimd}. First, they tested a bone conduction implant, a device that captures sound via an external microphone and transmits it wirelessly through the skull to the cochlea using an internal magnet \cite{Seo2021767}. Second, they tested nonmagnetic devices, such as pacemakers and \gls{icd} which are generally not significantly affected by magnetic fields because their enclosures are typically made of \gls{mri}-compatible metallic materials like titanium or zirconium.

The study showed that the bone conduction implant consistently exhibited the highest average torque (9.460 N $\cdot$ cm at 1.5 T and 14.618 N $\cdot$ cm at 3.0 T), while the pacemaker showed the lowest (0.255 N $\cdot$ cm at 1.5 T and 0.668 N $\cdot$ cm at 3.0 T). \gls{aimd} containing magnets, such as the cochlear and bone conduction implants, generally showed significantly higher torques than devices without magnets like \gls{icd} and pacemakers.

%-------------------------------------------------------------------------------------------------------------------
%-------Stoianovici 2024------------------TORQUE--------------------------------------------
%-------------------------------------------------------------------------------------------------------------------


Similarly, Stoianovici \textit{et al.} (2024) showcases a variation of the ASTM F2213 Torsional Spring measurement technique, referred as the 2-string method \cite{stoianovici2024}. This method address the difficulty of measuring very low magnetically induced torques on lightweight, slender medical devices, such as needles. This approach utilizes two strings to suspend the device by its ends and creates a torsional spring mechanism that is frictionless.

The needle rotates due to magnetic torque, the torque is calculated analytically based on the measured rotation angle and compared to the weight of the object. The study showed that the method successfully measured torques between $10^{-4} N \cdot m$ and 1.8$\times10^{-3}$ N$\cdot$m in the \gls{mri} environment. The maximum magnetically induced deflection angle remained below the 25$^\circ$ limit required by the ASTM F2213.

%RF and gradient fields

\subsection{Quantification of Magnetically Induced Heating}
\label{sec:rfpulses-heating}

Similarly, the ASTM F2182 guideline offers a standardized test procedure for detecting \gls{rf}-induced heating on or near passive implants in tissue-mimicking phantoms in order to quantify the heating effects of \gls{mri} scans on \gls{aimd}. Similar methods were used in earlier research by Aboyewa \textit{et al.} (2021) to examine the interaction of leads with the \gls{rf} field, assessing the energy absorbed and the ensuing heat generation \cite{aboyewa2021}.

The experiment was performed on a 3T Siemens Tim Trio\textsuperscript{\texttrademark}  \gls{mri}  scanner using a turbo spin echo pulse sequence. The pulse sequence was set to produce a whole-body Specific Absorption Rate (SAR) of approximately 2 W/kg and fluoroptic probes where used to measure the temperature rise. Several factors on induced heating were tested, including lead length (such as 6.33 cm, 12.65 cm, and 18.98 cm), lead type (bipolar vs. unipolar), orientation, and shape.

The study showed that lead length, type, and orientation significantly influence heating behavior while short leads ($<13$ cm) posed no significant thermal hazard during \gls{mri} at a specific absorption rate (SAR) of 2 W/kg, the maximum temperature rise was observed at the lead tips. For both lead types investigated, the heat induced increased with increasing lead length. Notably, bipolar leads showed slightly higher heating than unipolar leads and the straight lead configuration produced the highest maximum temperature rise \cite{aboyewa2021}.

%Significance
%Conclusion
The overall significance of these studies is in establishing reliable experimental and computational frameworks for quantifying \gls{mri}-induced mechanical effects, particularly where conventional ASTM approaches fall short. Stoianovici \textit{et al.} (2024) solves the technical challenge of measuring very low magnetically induced torques on slender, lightweight medical devices, such as needles, where conventional ASTM F2213 methods often become unreliable due to friction bias with their new methodology variant. Similarly, Seo \textit{et al.} (2021) validates the pulley methodology for \gls{aimd}. The modeling of clinically relevant alloys done by Woods \textit{et al.} (2021) shows experimental basis to rely on the listed alloys but requires further testing in terms of torque. While the modeling of the static magnetic field done by Ferreira \textit{et al.} (2017) provided a practical way to accurately map the magnetic field and calculate its spacial gradient. Finally, measurements by Aboyewa \textit{et al.} (2021) focused on the \gls{rf}-heating rather than mechanical effects. They employed the same \gls{mri} system and testing framework we utilize here, reinforcing the feasibility of our experimental setup and situating our work within established practices.
